\section{Pregunta N$^{\circ}$27\qquad Nombre}

\begin{frame}
	\begin{enumerate}\setcounter{enumi}{26}
		\item

		      La empresa de robótica Nobot produce tres robots de
		      juguete, modelos Dot, Lea y Liz.
		      Para fabricar cada robot de juguete, la empresa debe
		      utilizar las horas de mano de obra para codificar, el
		      montaje y la pintura.
		      Las cantidades para cada robot se muestran en la siguiente
		      tabla:

		      \begin{table}[ht!]
			      \centering
			      \begin{tabular}{c|c|c|c}
				      \hline
				               & Dot & Lea & Liz \\
				      \hline
				      Coding   & $4$ & $5$ & $1$ \\
				      Assembly & $7$ & $9$ & $2$ \\
				      Painting & $4$ & $2$ & $1$ \\
				      \hline
			      \end{tabular}
		      \end{table}

		      Hay $165$ horas de mano de obra disponibles para la
		      codificación, $295$ horas de mano de obra para el montaje
		      y 150 horas de mano de obra disponible para pintar cada
		      día.
		      ¿Cuántos de cada modelo se deben producir en un día, si se
		      utilizan todas las horas de mano de obra?
		      Utilice el método de eliminación gaussiana con pivote
		      parcial mediante la factorización $PA=LU$.
	\end{enumerate}

	\begin{solution}
		.
	\end{solution}
\end{frame}

\begin{frame}
	\begin{solution}
		.
	\end{solution}
\end{frame}